% P15.38: Hempel-Oppenheim 1948 Part IV, complete finite and measure-based routes.
\chapter{Systematic Power, Content and Logical Probability}
\label{sec:p15-systematic-power}

How much of a specified body of information can a theory
organize deductively? Hempel and Oppenheim propose a
measure of this task in Part IV of their 1948 paper
\cite[pp.~164--173]{HempelOppenheim1948Explanation}.
They first replace arbitrary sentence counts by a
fixed normal form, then relate the resulting measure
to logical probability. This chapter supplies both
constructions and their complete proofs.

Systematic power here is relative to a language,
data and a measure. It is not a test of a theory's
truth, an entropy measure or the explanation criterion
of \cref{def:p16-dn-potential-explanation}.
In this part of the source \(T\) may be any sentence,
including a singular or false sentence.

\section{Ground states and the content normal form}
\label{sec:p15-power-normal-forms}

\begin{hbdefinition}{def:p15-power-language}{The finite ground vocabulary}
Fix a first-order relational language with finitely
many constants and predicate symbols of finite arity,
without identity or nonconstant function symbols.
Its distinct ground atoms form a finite list
\(\mathsf{At}=\{A_1,\ldots,A_n\}\).
A singular sentence is a Boolean combination of
these atoms. The classical truth constants are
allowed, including for the degenerate case \(n=0\).
Quantifiers still range over nonempty structures;
the list of constants is not a domain-closure axiom.

Let \(\Omega=\{0,1\}^n\) be the formal ground
truth assignments. For \(v\in\Omega\), define
\[
 D_v=\bigwedge_{i=1}^n
       \begin{cases}A_i,&v_i=1,\\\neg A_i,&v_i=0,\end{cases}
 \qquad
 M_v=\bigvee_{i=1}^n
       \begin{cases}\neg A_i,&v_i=1,\\A_i,&v_i=0.\end{cases}
\]
The empty conjunction is \(\top\) and the empty
disjunction \(\bot\).
\(D_v\) is the source's maximal sentence or
state description; \(M_v\equiv\neg D_v\) is a
minimal sentence indexed by the assignment
it excludes.
For singular \(S\), put
\[
 [S]=\{v:v\models S\},\qquad
 \mathcal C(S)=\{M_v:v\notin[S]\}.
\]
The elements of \(\Omega\) are formal assignments,
not automatically equally plausible physical worlds.
\index{systematic power}\index{logical content!minimal sentences}
\end{hbdefinition}

\begin{hbtheorem}{thm:p15-power-normal-forms}{proved=annex:p15.power-normal-forms}{The complete maximal and minimal normal forms}
Every singular sentence satisfies
\[
               S\equiv\bigvee_{v\in[S]}D_v
                    \equiv\bigwedge_{M\in\mathcal C(S)}M .
\]
Both constituent sets are unique.
Distinct state descriptions are incompatible;
distinct minimal sentences have a valid disjunction.
A singular consequence of \(M_v\) is either
equivalent to \(M_v\) or valid. A singular
sentence entailing \(D_v\) is either equivalent
to \(D_v\) or contradictory.

There are \(2^n\) state descriptions and
\(2^n\) minimal sentences, and
\[
                   |[S]|+|\mathcal C(S)|=2^n .
\]
A valid sentence has empty content normal form;
a contradictory sentence has all minimal
sentences in that form. These statements
also hold for \(n=0\).
\end{hbtheorem}

\section{The original optimal-split construction}
\label{sec:p15-power-optimal-split}

\begin{hbdefinition}{def:p15-power-optimal-split}{Systematizing the minimal content of data}
Let \(K\) be a finite set of singular sentences,
let \(S=\bigwedge K\), and suppose \(S\) is
not valid. Write \(K'=\mathcal C(S)\).
For any sentence \(T\), consider the partitions
\[
              K'=K_1'\mathbin{\dot\cup}K_2'
       \quad\hbox{such that}\quad \{T\}\cup K_1'\entails M
                        \ \text{for every }M\in K_2'.
\]
The dot indicates a disjoint union. The source's
original systematic-power measure is
\[
 s_0(T,K)=\max\frac{|K_2'|}{|K'|}.
\]
The maximum exists because \(K'\) is finite and
nonempty. The partition with \(K_2'=\varnothing\)
is always available. Valid data are excluded
because their denominator is zero.
\end{hbdefinition}

\begin{hbtheorem}{thm:p15-power-optimal-split}{proved=annex:p15.power-optimal-split}{The unique optimal partition and source Theorem 8.2}
If a subfamily of distinct minimal sentences
is derivable from the remaining members
using \(T\), then every member of that
subfamily already follows from \(T\) alone.
Consequently the unique optimal second
part is
\[
 K_T=\{M\in\mathcal C(S):T\entails M\},
       \qquad s_0(T,K)=\frac{|K_T|}{|\mathcal C(S)|}.
\]
This holds even for a quantified \(T\),
and for an inconsistent \(T\).
It is the complete content of source
Theorems (8.2), (8.2a) and (8.2b).
\end{hbtheorem}

\begin{hbdefinition}{def:p15-power-uniform-measures}{Uniform range and content on singular sentences}
For singular \(U\), put
\[
 r_1(U)=\frac{|[U]|}{2^n},\qquad
 g_1(U)=\frac{|\mathcal C(U)|}{2^n}=1-r_1(U).
\]
Here content is the fraction of formal ground
assignments ruled out. It is not the number
of symbols, the number of premises or the
negative logarithm of a probability.
\end{hbdefinition}

\begin{hbtheorem}{thm:p15-power-uniform-formula}{proved=annex:p15.power-uniform-formula}{Representation invariance and the exact singular-theory formula}
For singular \(T,S\), with \(S\) not valid,
\[
                  s_0(T,\{S\})
                     =\frac{g_1(T\lor S)}{g_1(S)} .
\]
Replacing a finite data set by any logically
equivalent singular conjunction preserves
the measure. For a quantified \(T\), the
original optimal-split measure remains
defined, but this formula must use its
singular consequence projection
\[
              T^\flat=\bigwedge\{M_v:T\entails M_v\}
\]
in place of \(T\). The full truth of \(T\)
is not generally determined by its ground
assignment.
\end{hbtheorem}

\begin{hbtheorem}{thm:p15-power-source-examples}{proved=annex:p15.power-source-examples}{The source examples and all regrouped data}
In the language with one constant \(a\)
and unary \(P,Q\), let
\[
                 T=\forall x(P(x)\to Q(x)),
                 \qquad K=\{P(a),Q(a)\}.
\]
Then
\[
 K'=\{P(a)\lor Q(a),\ P(a)\lor\neg Q(a),\
                              \neg P(a)\lor Q(a)\},
 \qquad s_0(T,K)=\frac13 .
\]
Only the third minimal sentence follows
from \(T\) alone.

For the language with constants \(a,b,c,d\)
and unary predicates \(P,Q\), let
\[
\begin{split}
 K_1={}&\{P(a),Q(a),\neg P(b),\neg Q(b),Q(c),\neg P(d)\},\\
 K_2={}&\{P(a)\land Q(a),\neg P(b)\land\neg Q(b),Q(c),\neg P(d)\},\\
 K_3={}&\{P(a)\land\neg Q(b),Q(a)\land\neg P(b),Q(c),\neg P(d)\}.
\end{split}
\]
The fraction of listed sentences individually
derivable from the others using \(T\) is
respectively \(1/3,0,1/4\).
Their conjunctions are nevertheless equivalent,
and the content-based score is the same
for all three:
\[
                        s_0(T,K_j)=\frac{175}{252}
                                             =\frac{25}{36}.
\]
\end{hbtheorem}

\section{General range measures}
\label{sec:p15-power-range}

\begin{hbdefinition}{def:p15-power-range-measure}{A range measure on sentences}
A range measure \(r\) assigns a real number
to every sentence of the fixed language,
subject to:
\[
 0\le r(U)\le1,\qquad r(U)=1\ \text{when }U\text{ is valid},
\]
and
\[
 r(U\lor V)=r(U)+r(V)
          \quad\text{when }U\land V\text{ is contradictory}.
\]
These are the source's assumptions (9.7).
They impose finite additivity; no countable
continuity principle is included.
Define the corresponding content measure by
\[
                            g(U)=1-r(U).
\]
\end{hbdefinition}

\begin{hbtheorem}{thm:p15-power-range-laws}{proved=annex:p15.power-range-laws}{All range and content identities needed below}
The stated assumptions imply:
\[
\begin{gathered}
 r(\bot)=0,\qquad r(\neg U)=1-r(U),\qquad g(U)=r(\neg U),\\
 U\equiv V\ \Longrightarrow\ r(U)=r(V),\\
 U\entails V\ \Longrightarrow\ r(U)\le r(V),\\
 r(U\lor V)=r(U)+r(V)-r(U\land V).
\end{gathered}
\]
The function \(g\) is well defined, takes
values in \([0,1]\), has \(g(\bot)=1\), and
\[
 g(U\land V)=g(U)+g(V)
                 \quad\text{if }U\lor V\text{ is valid}.
\]
Thus it satisfies every source condition (9.8).
Examples of range measures include a probability
mixture of truth in finitely many fixed structures,
and, on the singular fragment, any weighted
distribution on \(\Omega\).
\end{hbtheorem}

\begin{hbdefinition}{def:p15-power-dual-measures}{Logical probability and generalized systematic power}
For a fixed range measure define
\[
 p_r(T,S)=\frac{r(T\land S)}{r(S)}
                       \quad\text{when }r(S)>0,
 \qquad
 s_r(T,S)=\frac{g(T\lor S)}{g(S)}
                       \quad\text{when }g(S)>0.
\]
The source writes these functions as \(p,s\).
They can have different domains: both
are defined for the same \(S\) exactly
when \(0<r(S)<1\).
\index{logical probability!range measure}
\index{systematic power!probability dual}
\end{hbdefinition}

\begin{hbtheorem}{thm:p15-power-probability-dual}{proved=annex:p15.power-probability-dual}{The complete complement relationship}
Whenever \(s_r(T,S)\) is defined,
\[
                    s_r(T,S)=p_r(\neg T,\neg S).
\]
Both sides have the same positive denominator.
In particular \(0\le p_r(T,S)\le1\) and
\(0\le s_r(T,S)\le1\) on their respective
domains. Under a fixed model probability,
the systematic score is the conditional
probability that \(T\) is false given that
\(S\) is false; it is not generally
the probability of \(T\) given \(S\).
\end{hbtheorem}

\begin{hbtheorem}{thm:p15-power-eight-identities}{proved=annex:p15.power-eight-identities}{All eight source probability and power identities}
For every displayed term that is defined,
\[
\begin{gathered}
 0\le p_r(T,S)\le1,\qquad 0\le s_r(T,S)\le1,\\
 p_r(\neg T,S)=1-p_r(T,S),\qquad
 s_r(\neg T,S)=1-s_r(T,S),\\[2pt]
 \begin{aligned}
 p_r(T_1\lor T_2,S)&=
 p_r(T_1,S)+p_r(T_2,S)-p_r(T_1\land T_2,S),\\
 s_r(T_1\land T_2,S)&=
 s_r(T_1,S)+s_r(T_2,S)-s_r(T_1\lor T_2,S),\\
 p_r(T_1\land T_2,S)&=p_r(T_1,S)\,p_r(T_2,T_1\land S),\\
 s_r(T_1\lor T_2,S)&=s_r(T_1,S)\,s_r(T_2,T_1\lor S).
 \end{aligned}
\end{gathered}
\]
The last probability identity requires
\(r(T_1\land S)>0\). If it is zero and
\(r(S)>0\), the left side is zero but
the second factor on the right is undefined.
The power identity has the corresponding
condition \(g(T_1\lor S)>0\); when that
value is zero and \(g(S)>0\), its left
side is zero and its second factor is
undefined. Neither case licenses multiplying
zero by an undefined conditional.
\end{hbtheorem}

\begin{hbdefinition}{def:p15-power-dual-syntax}{The dual of a Boolean argument}
For Boolean expressions in sentence metavariables
\(A_1,\ldots,A_m\), define \(D\) recursively:
\[
\begin{aligned}
 D(A_i)&=A_i,\qquad D(\neg B)=\neg D(B),\\
 D(B\land C)&=D(B)\lor D(C),\\
 D(B\lor C)&=D(B)\land D(C),\\
 D(\top)&=\bot,\qquad D(\bot)=\top.
\end{aligned}
\]
Expand implication as \(\neg B\lor C\)
before applying this operation.
Numerical addition, multiplication, constants
and order comparisons are left unchanged.
\end{hbdefinition}

\begin{hbtheorem}{thm:p15-power-dualism}{proved=annex:p15.power-dualism}{A precise form of source Theorem 9.13}
Every universally valid equality or inequality
formed from defined \(p_r(B,C)\) terms,
Boolean sentence arguments and numerical
operations yields a universally valid
statement when each \(p_r(B,C)\) is
replaced by \(s_r(D(B),D(C))\).
The converse transformation holds as well.
The definedness conditions transform along
with the terms.

Boolean truth constants must be exchanged
as in the definition. If a statement
additionally uses an entailment hypothesis,
the complemented substitution reverses
its direction:
\[
 B(\neg A_1,\ldots,\neg A_m)\entails
 C(\neg A_1,\ldots,\neg A_m)
 \quad\Longleftrightarrow\quad
 D(C)\entails D(B).
\]
An informal symbol exchange without these
conditions is not the theorem.
\end{hbtheorem}

\section{What the numerical score does and does not establish}
\label{sec:p15-power-interpretation}

\begin{hbtheorem}{thm:p15-power-strength-boundary}{proved=annex:p15.power-strength-boundary}{Strength, self-systematization and zero-weight countermodels}
For \(g(S)>0\):
\[
 s_r(\top,S)=0,\qquad s_r(\bot,S)=1,\qquad
 s_r(S,S)=1.
\]
If \(T\entails U\), then
\(s_r(T,S)\ge s_r(U,S)\).
If \(T\entails S\), then \(s_r(T,S)=1\).
The converse need not hold when countermodels
have zero measure.

On a finite ground space with every assignment
given positive weight and singular \(T,S\),
\[
 s_r(T,S)=1\quad\Longleftrightarrow\quad T\entails S.
\]
When \(r(S)>0\), full support also gives
\(p_r(T,S)=1\) exactly when \(S\entails T\).
Without full support either converse can fail.
A maximal power score
therefore does not by itself establish truth,
independent explanation or scientific merit.
\end{hbtheorem}

\begin{hbtheorem}{thm:p15-power-weight-choice}{proved=annex:p15.power-weight-choice}{Two explicit range choices and their different learning behavior}
For two singular atoms \(A,B\), a uniform
distribution on their four assignments gives
\[
                  p_r(B,A)=s_r(B,A)=\frac12.
\]
Instead, choose a latent parameter
\(\theta\in\{1/4,3/4\}\), each value with
probability \(1/2\), and let \(A,B\) be
independent Bernoulli outcomes conditional
on \(\theta\). The induced weights are
\[
 (w_{00},w_{01},w_{10},w_{11})
                         =(5/16,3/16,3/16,5/16),
\]
and
\[
                  p_r(B,A)=s_r(B,A)=\frac58.
\]
Both choices satisfy the range axioms.
The second is an explicit illustration
of different weights, not a claim to
reproduce Carnap's particular measure.
\end{hbtheorem}

\begin{hbtheorem}{thm:p15-power-language-choice}{proved=annex:p15.power-language-choice}{Fixed-language invariance is not invariance under every new vocabulary}
For singular \(T,S\), adding fresh, unconstrained
ground atoms and retaining the uniform measure
preserves \(s_0(T,\{S\})\).
Changing the primitive representation can
change the uniform score: \(T=A,S=B\)
has score \(1/2\) on two independent formal
atoms. Replacing \(A\) by \(C\lor D\)
and using the uniform measure on \(B,C,D\)
gives score \(1/4\).

For the quantified source theory
\(T=\forall x(P(x)\to Q(x))\) and
the same data \(S=P(a)\land Q(a)\),
the original optimal-split score is \(1/3\)
with the sole name \(a\) and \(1/2\)
with the names \(a,b\).
The new ground instances of a quantified
theory need not be unconstrained.
\end{hbtheorem}

\begin{hbtheorem}{thm:p15-power-finite-procedure}{proved=annex:p15.power-finite-procedure}{A complete finite calculation and its first-order boundary}
For singular \(T,S\), enumerate the \(2^n\)
assignments and count those falsifying \(S\)
and those falsifying both \(T,S\).
If the first count is positive, their ratio
is exactly the uniform systematic score.
With specified nonnegative rational probability weights,
sum weights instead of counts; the result
is exact whenever the denominator is positive.

For quantified \(T\), computing the original
optimal-split score instead asks, for each
assignment \(v\) falsifying \(S\), whether
\(T\land D_v\) has a first-order model.
An absent model contributes the minimal
sentence \(M_v\) to the numerator.
This is a finite number of first-order
queries, not a general terminating procedure
for deciding those queries. A single ground
assignment can have both a model of \(T\)
and a model of \(\neg T\) as extensions.
\end{hbtheorem}

The source concludes by leaving two application
questions open: how to specify suitable qualitative
primitives, and how to justify a range measure
among the many formal possibilities.
Its references to Carnap and to Helmer,
Hempel and Oppenheim name distinct confirmation
constructions; their detailed measures and
limiting extensions are not supplied by
the general ratio formulas alone.
The definitions here make a specified
calculation reproducible while keeping
those further choices explicit.
